\documentclass{article}

\usepackage{amsfonts}

\begin{document}

\textbf{Chapter 2
Qualitative properties and predictability
of evolutions}

\hfill

In the previous chapter we already met various types of evolutions, including
stationary and periodic ones. In the present chapter we investigate these evolutions
more in detail and we also consider more general types. Here qualitative properties
are our main interest. The precise meaning of ‘qualitative’ in general is hard to give;
to some extent it is the contrary or opposite of ‘quantitative’. Roughly one could say
that qualitative properties can be described in words (or integers), whereas quantitative
properties usually are described in terms of (real) numbers.

\hfill

It turns out that for the various types of evolutions, it is not always equally easy to
predict the future course of the evolution. At first sight this may look strange: once
we know the initial state and the evolution operator, we can just predict the future
of the evolution in a unique way. However, in general, we only know the initial
state approximately. Moreover, often the evolution operator is not known at all or
only approximately. We ask the question of predictability mainly in the following
form. Suppose the evolution operator is unknown, but that we do know a (long)
piece of the evolution; is it then possible to predict the future course of the evolution
at least partially? It turns out that for the various kinds of evolution, the answer to
this question strongly differs. For periodic, in particular for stationary, evolutions,
the predictability is not a problem. Indeed, based on the observed regularity we
predict that this will also continue in the future. As we show, the determination of
the exact period still can be a problem. This situation occurs, to an even stronger
degree, with quasi-periodic solutions. However, it turns out that for the so-called
chaotic evolutions an essentially worse predictability holds. In Chapter 6 we return
to predictability problems in a somewhat different context.

\hfill

\textbf{2.1 Stationary and periodic evolutions}

\hfill

Stationary and periodic evolutions have been already widely discussed in the previous
chapter.

\hfill

\textbf{Definition 2.1 (Stationary and periodic evolutions).} Let $(M, T, \Phi)$; be a dynamical
system with state space M, time set T, and evolution operator $\Phi: M \times T \to M$.

\hfill

We say that $x \in M$ is periodic (or stationary) if there exists a positive t 2 T such
that $\Phi(x,t) = x$.

\begin{enumerate}
\item Moreover, x is stationary whenever for all $t \in T$ one has $\Phi(x,t)=x$
\item If $\Phi(x, \tau)=x$, we call $\tau$ a period of x:
\end{enumerate}

The periodic or stationary evolution of x as a set is given by

\[
\{ \Phi(x,s) | s \in T|\} \subset M
\]

We already noted in the example of the doubling map (see 1.3.8) that each
periodic point has many periods: if $\tau$ is a period, then so are all integer multiples
of $\tau$. As before we introduce the notion of prime period for periodic, nonstationary,
evolutions. This is the smallest positive period. This definition appears to be without
problems, but this is not entirely true. In the case where $T = \mathbb{R}$, such a prime period
does not necessarily exist. First of all, this is the case for stationary evolutions,
which can be excluded by just considering proper, that is, nonstationary, periodic
evolutions. However, even then there are pathological examples; see Exercise 2.4.
In the sequel we assume that prime periods of nonstationary periodic evolutions
will always be well defined. It can be shown that this assumption is valid in all
cases where M is a metric (or metrisable) space and where the evolution operator,
considered as a map $\Phi \times T \to M$, is continuous. If such a prime period equals $\tau$, it is easy to show that the set of periods is given by $\tau\mathbb{Z} \bigcap T$ 

\hfill

\textbf{2.1.1 Predictability of periodic and stationary motions}

\hfill

The prediction of future behaviour in the presence of a stationary evolution is
completely trivial: the state does not change and the only reasonable prediction is
that this will remain so. In the case of periodic evolution, the prediction problem
does not appear much more exciting. In the case where $T = \mathbb{Z}$ or $T = \mathbb{Z}_{+}$; this is
indeed so. Once the period has been determined and the motion has been precisely
observed over one period, we are done. We show, however, that in the case where 
$T = \mathbb{R}$ or $T = \mathbb{R}_{+}$, when the period is a real number, the situation is less clear.We
illustrate this when dealing with the prediction of the yearly motion of the seasons,
or equivalently, the yearly motion of the Sun with respect to the vernal equinox (the
beginning of Spring). The precise prediction of this cycle over a longer period is
done with a ‘calendar’. In order to make such a calendar, it is necessary to know the
length of the year very accurately and it is exactly here where the problem occurs.
Indeed, this period, measured in days, or months, is not an integer and we cannot
expect to do better than find an approximation for this. By the ‘calendar’ here we
should understand the set of rules according to which this calendar is made, where
one should think of the leap year arrangements. Instead of ‘calendar’ one also speaks
of time tabulation.

\hfill

Remark (Historical digression).We give some history. In the year 45 BC the Roman
dictator Julius Cæsar 2 installed the calendar named after him, the Julian Calendar.
The year was thought to consist of 365.25 days. On the calendar this was implemented
by adding a leap day every four years. However, the real year is a bit shorter
and slowly the seasons went out of phase with the calendar. Finally Pope Gregorius
XIII 3 decided to intervene. First, to ‘reset’ the calendar with nature, he decided that
in the year 1582 (just at this one occasion) immediately after October 4, October
15 would follow. This means that the predictions based on the Julian system in this
period of more than 1600 years already show an error of 10 days. Apart from this
unique intervention and in order to prevent such shifts in the future, the leap year
conventionwas changed as follows. In the future, year numbers divisible by 100, but
not by 400, no longer will be leap years. This led to the present, Gregorian, calendar.
Compare [158].

\hfill

Of course we also understand that this intervention is still based on an approximation:
real infallibility cannot be expected here! In fact, a correction of 3 days per
400 years amounts to 12 days per 1600 years, whereas only 10 days were necessary.

\hfill

Our conclusion from this example is that when predicting the course of a periodic
motion, we expect an error that in the beginning is proportional to the interval of
prediction, that is, the time difference between the present moment of prediction
and the future instant of state we predict. When the prediction errors increase, a
certain saturation occurs: when predicting seasons the error cannot exceed one half
year. Moreover, when the period is known to a greater precision (by longer and/or
more precise observations), the growth of the errors can be made smaller.

\hfill

As said earlier, the problem of determination of the exact period only occurs in
the case of time sets $\mathbb{R}$ or $\mathbb{R}_{+}$. In the cases where the time sets are $\mathbb{Z}$ or $\mathbb{Z}_{+}$ the
periods are integers, which by sufficiently accurate observation can be determined
with absolute precision. These periods then can be described by a counting word;
they can be viewed as qualitative properties. In cases where the time set is $\mathbb{R}$ or $\mathbb{R}_{+}$,
however, the period is a quantitative property of the evolution.

\hfill

On the other hand the difference is not that sharp. If the period is an integer,
but very large, then it can be practically impossible to determine this by observation.
As we show when discussing so-called multiperiodic evolutions, it may even
be impossible to determine by observation whether we are dealing with a periodic
evolution.

\hfill


Also, regarding the more complicated evolutions dealt with below, we ask ourselves
how predictable these are and which errors we can expect as a function of the
prediction interval. We then merely consider predictions based on a (sufficiently)
long observation of the evolution.

\hfill

\textbf{2.1.2 Asymptotically and eventually periodic evolutions}

\hfill

Returning to the general context of dynamical systems we now give the following
definition.

\hfill

\textbf{Definition 2.2 (Asymptotic periodicity and stationarity).} Let $(M, T, \Phi)$ be a
dynamical system, where M is a metric space and $\Phi: M \times T \to M$ is continuous.
We call the evolution $x(t) = \Phi(x,t)$ asymptotically periodic (or stationary) if for $t \to \infty$ the point x(t) tends to a periodic (or stationary) evolution $\tilde{x}(t)$, in the sense
that the distance between x(t) and the set $\{ \tilde{x}(s) | s \in T \}$ tends to zero as $t \to \infty$.

\hfill

Remarks.

\begin{itemize}
\item This definition can be extended to the case where M only has a topology. In that
case one may proceed as follows.We say that the evolution x(t) is asymptotically
periodic if there exists a periodic orbit $\gamma = \{ \tilde{x}(s) | s \in T \}$ such that for any
neighbourhood U of $\gamma$ there exists $T_U \in \mathbb{R}$ such that for $t > T_U$ one has that
$x(t) \in U$.

\hfill

We speak of asymptotic periodicity in the strict sense if there exists a periodic
orbit $\gamma = \{ \tilde{x}(s) | s \in T \}$ of period $\tau$ and a constant c, such that

\[
lim_{\mathbb{N} \ni \to \infty} x(t + n \tau)) = \tilde{x}(t+c)
\]

In that case either the number $c = \tau$ modulo $\mathbb{Z}$ or the point $\tilde{x}(c)$ can be interpreted
as the phase associated with x(t).

\hfill

Note that in the case without topology, the concept of ‘asymptotic’ no longer
makes sense, because the concept of limit fails to exist. The eventually periodic
or stationary evolutions, discussed later in this section, are an exception to this.

\item Assuming continuity of $\Phi$, it follows that the set $\{ \tilde{x}(s) | s \in T \}$ corresponding
to the periodic evolution $\tilde{x}(s)$ is compact. Indeed, in the case where $T = \mathbb{Z}$ or $v_{+}$ the set is even finite, whereas in the cases where $T = \mathbb{R}$ or $\mathbb{R}_{+}$ the set $\{ \tilde{x}(s) | s \in T \}$ is the continuous image of the circle. The reader can check this
easily.

\end{itemize}

We have seen examples of asymptotically periodic or asymptotically stationary
evolutions several times in the previous chapter: for the free pendulum with damping
(see $\S$1.1.2), every nonstationary evolution is asymptotically stationary. For the
damped pendulum with forcing (see $\S$1.1.2), and in the Van der Pol system (see $\S$1.3.1)we saw asymptotically periodic evolutions.When discussing the H´enon family
(in $\S$1.3.2) we saw another type of stationary evolution, given by a saddle point.
For such a saddle point p we defined stable and unstable separatrices (see (1.16)).
From this it may be clear that the stable separatrix $W^s(p)$ exactly consists of points
that are asymptotically stationary and of which the limit point is the saddle point p:

\hfill

The definition of asymptotically periodic or asymptotically stationary applies for
any of the time sets we are using. In the case of time set $\mathbb{Z}_{+}$ another special type
occurs.

\hfill

\textbf{Definition 2.3 (Eventual periodicity and stationarity).} Let $(M, \mathbb{Z}_{+}, \Phi)$ be a
dynamical system with an evolution $x(n) = \Phi(x,n)$. We say that x(n) is eventually
periodic (or eventually stationary) if for certain N >0 and k > 0,

\[
x(n) = x(n+k) for all n > N.
\]

So this is an evolution that after some time is exactly periodic (and not only in a
small neighbourhood of a periodic evolution). Note that in this case neither distance
function nor topology plays any role.

\hfill

This type of evolution is abundantly present in the system generated by the doubling
map $x \mapsto 2x \mod \mathbb{Z}$, defined on [0, 1); see $\S$1.3.8. Here 0 is the only stationary
point. However, any rational point p/q with q a power of 2 is eventually stationary.
Compare Exercise 2.3. Similarly, any rational point is eventually periodic. Also for
the logistic system (see $\S$1.3.3), such eventually periodic evolutions occur often.

\hfill

As said before, such evolutions do not occur in the case of the time sets $\mathbb{R}$ and \mathbb{Z}.
Indeed, in these cases, the past can be uniquely reconstructed from the present state,
which clearly contradicts the existence of an evolution that ‘only later’ becomes
periodic. (Once the evolution has become periodic, we cannot detect how long ago
this has happened.)

\hfill

Regarding the predictability of these evolutions the same remarks hold as in
the case of periodic and stationary evolutions. As said earlier, when discussing
predictability, we always assume that the evolution to be predicted has been observed
for quite some time. This means that it is no further restriction to assume
that we already are in the stage that no observable differences exist between the
asymptotically periodic (or stationary) evolutions and the corresponding periodic
(or stationary) limit evolution.

\hfill

\textbf{2.2 Multi- and quasi-periodic evolutions}

\hfill

In the previous chapter we already met with evolutions where more than one period
(or frequency) plays a role. In particular this was the case for the free pendulum
with forcing; see $\S$1.1.2. Also in daily life we observe such motions. For instance,
the motion of the Sun as seen from the Earth knows two periodicities: the daily
period, of sunrise and sunset, and the yearly one that determines the change of the
seasons.

\hfill

We describe these types of evolution as a generalised type of periodic motions,
starting out from our dynamical system $(M, T, \Phi)$. We assume that M is a differentiable
manifold and that $\Phi: M \times T \to M$ is a differentiable map. For periodic
evolutions, the whole evolution $\{ x(t) \ t \in T \}$ as a subset of the state spaceM consists
either of a finite set of points (in cases where $T = \mathbb{Z}_{+}$ or $\mathbb{Z}$ or of a closed curve
(in the case where $T = \mathbb{R}$ or $\mathbb{R}_{+}$. Such a finite set of points we can view as $\mathbb{Z}/s\mathbb{Z}$,
where s is the (prime-) period: the closed curve can be viewed as the differentiable

\end{document}
